% Standalone insertion draft. The original manuscript has not been modified.
% Quantitative placeholders are filled only after the audited formal run.

\subsection{Ablation Analysis of the ART-MAPPO Components}
\label{subsec:art_mappo_ablation}

To isolate the contribution of the three architectural mechanisms requested
by the reviewer, we compare the complete ART-MAPPO against three
single-component ablations: without trust-modulated exploration, without the
GRU temporal encoder, and without the attention--residual backbone. The
environment, reward, action limits, target assignment, optimizer, rollout
budget, and random seeds are identical. In the last ablation, attention and
residual processing are replaced by a parameter-matched multilayer
perceptron, so that the comparison is not explained by a smaller network.
Three independent training seeds are used. After training, guided exploration
is disabled and the selected checkpoints are tested on the same 100 paired
Monte Carlo episodes in each engagement case.

\begin{figure}[!t]
\centering
\includegraphics[width=\columnwidth]{ABLATED_TRAINING_CASE1_FIGURE}
\caption{Component ablation in Case~1 at the training level: episode return,
critic loss, and policy entropy. Solid lines denote the three-seed mean and
shaded regions denote one standard deviation.}
\label{fig:ablation_training_case1}
\end{figure}

\begin{figure}[!t]
\centering
\includegraphics[width=\columnwidth]{ABLATED_TRAINING_CASE2_FIGURE}
\caption{Component ablation in Case~2 at the training level: episode return,
critic loss, and policy entropy. Solid lines denote the three-seed mean and
shaded regions denote one standard deviation.}
\label{fig:ablation_training_case2}
\end{figure}

The trust-aware mechanism is interpreted first from its role in the proposed
method: it is a training-time memory that changes the probability of using
tactical guidance, probe, and random actions according to standardized
episodic returns. It is disabled at deployment and is therefore expected to
affect the learned policy indirectly through exploration coverage and the
entropy trajectory, rather than to act as an additional test-time controller.
The no-trust comparison consequently focuses on reward acquisition,
inter-seed stability, and policy-entropy decay; the frozen-policy Monte Carlo
result is reported as the downstream consequence of that changed training
distribution.

The GRU and attention--residual backbone have roles during both optimization
and deployment. The GRU preserves engagement history and can infer temporal
trends from delayed and noisy line-of-sight observations. Its independent
effect is therefore evaluated through both training stability and the
frozen-policy interception/coordination rates. The attention mechanism
learns interactions among physically meaningful observation channels, while
the residual identity paths preserve primitive kinematic information and
improve optimization conditioning. The capacity-matched backbone ablation
tests whether these structured operations improve critic fitting and final
guidance quality beyond what can be obtained from a plain network of nearly
the same size.

\begin{figure}[!t]
\centering
\subfloat[Case~1]{\includegraphics[width=\columnwidth]{ABLATED_TEST_CASE1_FIGURE}}\\
\subfloat[Case~2]{\includegraphics[width=\columnwidth]{ABLATED_TEST_CASE2_FIGURE}}
\caption{Frozen-policy component ablation over 100 paired Monte Carlo
episodes per method and case. The bars report interception and cooperative
success with 95\% Wilson confidence intervals; the boxplots use the
unmodified successful-episode samples.}
\label{fig:ablation_monte_carlo}
\end{figure}

% RESULT_PARAGRAPHS_TO_BE_FILLED_FROM_CSV
